% GENERATED FILE. Edit canonical lessons, then run npm run generate:books.
% canonical-source-hash: fnv1a64:408d7006a5474323
% canonical-lessons: RU-C73-below, RU-C73-left, RU-C73-right, RU-C73-alot, RU-C73-alittle

\chapter{Below, Left, Right, A Lot, A Little}
\label{ch:ru-below-left-right-a-lot-a-little}
\hlchaptermodality{73}

\begin{chapteropening}
\textbf{By the end of this chapter:} \emph{I can say below, on the left, on the right, a lot and a little.}

Complete the last lesson of chapter 73: I can say below, on the left, on the right, a lot and a little.
\end{chapteropening}

\section[\texorpdfstring{\ru{внизу} (vnizú)}{внизу (vnizú)}]{\ru{внизу} (vnizú) — below, downstairs}

\label{lesson:RU-C73-below}



\begin{quote}
Before the new one: say the Russian for near, then the Russian for up above.
\end{quote}

\subsection*{You'll want to know: \ru{внизу}}
\textbf{\ru{внизу}} (\emph{vnizú}) — \textquotedblleft{}below, downstairs\textquotedblright{}.

\begin{cousinweb}
\textbf{\ru{внизу}} is \emph{\ru{в}}, \textquotedblleft{}in\textquotedblright{}, with \emph{\ru{низу}}, a form of \emph{\ru{низ}}, \textquotedblleft{}the bottom\textquotedblright{}. Vasmer explains \emph{\ru{низ}} as an extension of the PIE \emph{*ni}, \textquotedblleft{}down\textquotedblright{}, in Sanskrit \emph{ni-}, \textquotedblleft{}down\textquotedblright{}, and Old High German \emph{nidar}, \textquotedblleft{}down\textquotedblright{}: English \emph{nether} is a cousin.
\end{cousinweb}

\begin{grammarlens}[title={What you've built}]
One more word for this chapter.
\end{grammarlens}

\subsection*{Your turn}
\begin{itemize}
\raggedright
  \item \emph{Say it:} \emph{vnizú}
  \item \emph{Say it:} \emph{vnizú}, once more
  \item \emph{Say it:} say \emph{vnizú}
  \item \emph{Recall:} say the Russian for near, then the Russian for up above, then say \emph{vnizú} again
\end{itemize}

\begin{tcolorbox}[breakable,colback=teal!4,colframe=teal!35!black,title={Before you move on}]
What does \ru{внизу} mean? (\textquotedblleft{}Below, downstairs\textquotedblright{}.) Say it once more.
\end{tcolorbox}
\section[\texorpdfstring{\ru{слева} (sléva)}{слева (sléva)}]{\ru{слева} (sléva) — on the left}

\label{lesson:RU-C73-left}



\begin{quote}
Before the new one: say the Russian for up above, then the Russian for below.
\end{quote}

\subsection*{You'll want to know: \ru{слева}}
\textbf{\ru{слева}} (\emph{sléva}) — \textquotedblleft{}on the left\textquotedblright{}.

\begin{cousinweb}
\textbf{\ru{слева}} is \emph{\ru{с}}, \textquotedblleft{}from\textquotedblright{}, with \emph{\ru{лева}}, a form of \emph{\ru{левый}}, \textquotedblleft{}left\textquotedblright{}: \textquotedblleft{}from the left\textquotedblright{}. \emph{\ru{левый}} $\leftarrow$ Proto-Slavic \emph{*lěvŭ}, related to Latin \emph{laevus} and Greek \emph{laiós}, both \textquotedblleft{}left\textquotedblright{} (Vasmer).
\end{cousinweb}

\begin{grammarlens}[title={What you've built}]
One more word for this chapter.
\end{grammarlens}

\subsection*{Your turn}
\begin{itemize}
\raggedright
  \item \emph{Say it:} \emph{sléva}
  \item \emph{Say it:} \emph{sléva}, once more
  \item \emph{Say it:} say \emph{sléva}
  \item \emph{Recall:} say the Russian for up above, then the Russian for below, then say \emph{sléva} again
\end{itemize}

\begin{tcolorbox}[breakable,colback=teal!4,colframe=teal!35!black,title={Before you move on}]
What does \ru{слева} mean? (\textquotedblleft{}On the left\textquotedblright{}.) Say it once more.
\end{tcolorbox}
\section[\texorpdfstring{\ru{справа} (správa)}{справа (správa)}]{\ru{справа} (správa) — on the right}

\label{lesson:RU-C73-right}



\begin{quote}
Before the new one: say the Russian for below, then the Russian for on the left.
\end{quote}

\subsection*{You'll want to know: \ru{справа}}
\textbf{\ru{справа}} (\emph{správa}) — \textquotedblleft{}on the right\textquotedblright{}.

\begin{cousinweb}
\textbf{\ru{справа}} is \emph{\ru{с}}, \textquotedblleft{}from\textquotedblright{}, with \emph{\ru{права}}, a form of \emph{\ru{правый}}, \textquotedblleft{}right\textquotedblright{}, first \textquotedblleft{}straight, righteous\textquotedblright{}. Vasmer derives \emph{\ru{правый}}, probably, from \emph{*prō-vos}, \textquotedblleft{}forward\textquotedblright{}, related to Latin \emph{probus}, \textquotedblleft{}upright, good\textquotedblright{}.
\end{cousinweb}

\begin{grammarlens}[title={What you've built}]
One more word for this chapter.
\end{grammarlens}

\subsection*{Your turn}
\begin{itemize}
\raggedright
  \item \emph{Say it:} \emph{správa}
  \item \emph{Say it:} \emph{správa}, once more
  \item \emph{Say it:} say \emph{správa}
  \item \emph{Recall:} say the Russian for below, then the Russian for on the left, then say \emph{správa} again
\end{itemize}

\begin{tcolorbox}[breakable,colback=teal!4,colframe=teal!35!black,title={Before you move on}]
What does \ru{справа} mean? (\textquotedblleft{}On the right\textquotedblright{}.) Say it once more.
\end{tcolorbox}
\section[\texorpdfstring{\ru{много} (mnógo)}{много (mnógo)}]{\ru{много} (mnógo) — a lot, many}

\label{lesson:RU-C73-alot}



\begin{quote}
Before the new one: say the Russian for on the left, then the Russian for on the right.
\end{quote}

\subsection*{You'll want to know: \ru{много}}
\textbf{\ru{много}} (\emph{mnógo}) — \textquotedblleft{}a lot, many\textquotedblright{}.

\begin{cousinweb}
\textbf{\ru{много}} $\leftarrow$ Proto-Slavic \emph{*mŭnogŭ}, \textquotedblleft{}much, many\textquotedblright{} — inherited. Vasmer relates it to Gothic \emph{manags}, \textquotedblleft{}many\textquotedblright{}, and Old Irish \emph{menicc}, \textquotedblleft{}frequent\textquotedblright{}, and finds no grounds for a Germanic loan: English \emph{many} is a cousin.
\end{cousinweb}

\begin{grammarlens}[title={What you've built}]
One more word for this chapter.
\end{grammarlens}

\subsection*{Your turn}
\begin{itemize}
\raggedright
  \item \emph{Say it:} \emph{mnógo}
  \item \emph{Say it:} \emph{mnógo}, once more
  \item \emph{Say it:} say \emph{mnógo}
  \item \emph{Recall:} say the Russian for on the left, then the Russian for on the right, then say \emph{mnógo} again
\end{itemize}

\begin{tcolorbox}[breakable,colback=teal!4,colframe=teal!35!black,title={Before you move on}]
What does \ru{много} mean? (\textquotedblleft{}A lot, many\textquotedblright{}.) Say it once more.
\end{tcolorbox}
\section[\texorpdfstring{\ru{мало} (málo)}{мало (málo)}]{\ru{мало} (málo) — a little, few}

\label{lesson:RU-C73-alittle}



\begin{quote}
Before the new one: say the Russian for on the right, then the Russian for a lot.
\end{quote}

\subsection*{You'll want to know: \ru{мало}}
\textbf{\ru{мало}} (\emph{málo}) — \textquotedblleft{}a little, few\textquotedblright{}.

\begin{cousinweb}
\textbf{\ru{мало}} is \emph{\ru{малый}}, \textquotedblleft{}small\textquotedblright{}, made into an adverb with \emph{-\ru{о}}. \emph{\ru{малый}} is a cousin, Vasmer says, of Gothic \emph{smals}, \textquotedblleft{}small\textquotedblright{}, and so of English \emph{small}.
\end{cousinweb}

\begin{grammarlens}[title={What you've built}]
One more word for this chapter.
\end{grammarlens}

\subsection*{Your turn}
\begin{itemize}
\raggedright
  \item \emph{Say it:} \emph{málo}
  \item \emph{Say it:} \emph{málo}, once more
  \item \emph{Say it:} say \emph{málo}
  \item \emph{Recall:} say the Russian for on the right, then the Russian for a lot, then say \emph{málo} again
\end{itemize}

\begin{tcolorbox}[breakable,colback=teal!4,colframe=teal!35!black,title={Before you move on}]
What does \ru{мало} mean? (\textquotedblleft{}A little, few\textquotedblright{}.) Say it once more.
\end{tcolorbox}
